% Response to the Google Deep Review — Round 3
% For "Trustworthy Agentic Automation for Intent-Based Self-Healing in
% Software-Defined Networks: Three Deterministic Gates around a Bounded Agent"
%
% Compile with:  tectonic REVISION_REPORT_ROUND3.tex
\documentclass[11pt,a4paper]{article}
\usepackage[margin=1.9cm]{geometry}
\usepackage{mathptmx}
\usepackage{amsmath,amssymb}
\usepackage{booktabs,longtable,array,colortbl}
\usepackage{enumitem}
\usepackage{xcolor}
\usepackage{fancyhdr,titlesec}
\usepackage[hidelinks]{hyperref}
\usepackage{microtype}
\emergencystretch=2.5em
\sloppy

\definecolor{ink}{HTML}{1F2A44}
\definecolor{accent}{HTML}{0F5EA8}
\definecolor{quote}{HTML}{5A6B7C}

\hypersetup{
  pdftitle={Response to the Google Deep Review --- Round 3},
  pdfauthor={The author},
  colorlinks=true, urlcolor=accent, linkcolor=accent, citecolor=accent}

\titleformat{\section}{\Large\bfseries\color{ink}}{\thesection}{1em}{}
\titleformat{\subsection}{\large\bfseries\color{ink}}{\thesubsection}{0.8em}{}
\setlist{nosep,topsep=3pt}

\pagestyle{fancy}
\fancyhead[L]{\small\color{quote}Response to the Google Deep Review --- Round 3}
\fancyhead[R]{\small\color{quote}\thepage}
\fancyfoot{}

% Numbers quoted from the boundary probe (final_submission/data/boundary_probe.csv).
\newcommand{\ProbeTotal}{20}
\newcommand{\ProbeApproved}{13}
\newcommand{\ProbeRejected}{7}
\newcommand{\ProbeTelApproved}{8}
\newcommand{\ProbeTelN}{8}
\newcommand{\ProbeTelCross}{0}
\newcommand{\ProbeTopoN}{8}
\newcommand{\ProbeTopoApproved}{5}
\newcommand{\ProbeTopoCross}{2}
\newcommand{\ProbeTopoFaultyApproved}{0}
\newcommand{\ProbeCraftedN}{4}
\newcommand{\ProbeCraftedApproved}{0}
\newcommand{\ProbeStrong}{3}
\newcommand{\ProbeStrongApproved}{0}
\newcommand{\PromptCases}{6}

\newenvironment{reviewer}{\par\medskip\noindent\begin{minipage}{\linewidth}\small\itshape\color{quote}\textbf{Reviewer:}\par}{\end{minipage}\medskip\par}
\newenvironment{response}{\par\noindent\begin{minipage}{\linewidth}\small\textbf{Response.}\par}{\end{minipage}\medskip\par}

\begin{document}

\begin{titlepage}
\vspace*{2.8cm}
{\color{ink}\Huge\bfseries Response to the Google Deep Review\\[0.3em]\bigskip}
{\color{ink}\Large Round 3 Report\bigskip}
{\color{ink}\large\bfseries
\emph{Trustworthy Agentic Automation for Intent-Based Self-Healing in
Software-Defined Networks: Three Deterministic Gates around a Bounded Agent}\\[0.6em]}
{\large Revised submission, November 2026}

\vspace{1.2cm}
\begin{tabular}{@{}ll@{}}
\toprule
\textbf{Round 3 review} & \texttt{final\_submission/GA\_DEEP\_REVIEW\_ROUND3.pdf} \\
\textbf{Review source} & \texttt{GA\_DEEP\_REVIEW\_ROUND3.md} (byte-identical to the pasted text) \\
\textbf{Revised paper} & \texttt{final\_submission/paper/paper.pdf} \\
\textbf{Round 1 / 2 letters} & \texttt{REVISION\_REPORT.pdf}, \texttt{REVISION\_REPORT\_ROUND2.pdf} \\
\textbf{One command} & \texttt{bash final\_submission/scripts/reproduce.sh} \\
\bottomrule
\end{tabular}

\vspace{1.5cm}
\noindent\textbf{A note on how to read this document.}

Every Round~3 comment is reproduced in the reviewer's words (lightly
trimmed, never paraphrased beyond recognition) and followed by a numbered
response stating what changed and where to check it. The review's own
structure is kept. Two of the responses are longer than the rest, and they
are the two we would most like a re-reader to look at: the review found
that our own adversarial probe of the request screen and of the
clarification branch was missing two behaviours, and both were fixed with
regression tests before this letter was written (\S2.4 and \S2.6). Where we
did not change the artifact, we say so and why; where we changed only the
paper, we say that too.

\vfill
\end{titlepage}

\section{Summary of changes, and what the review actually found}

The Round~3 review confirmed the artifact's structure and its honesty---it
verifies the internal arithmetic of Tables~II--VIII against Eq.~(2) and
finds it consistent---and then identified five substantive gaps.

\begin{enumerate}
  \item \textbf{The framing of the contribution.} The review is right that
    the paper's related work conflated three different families of prior
    systems, and that our claim (``none gates the action on a post-change
    measured outcome verdict'') was too broad as written. Section~II-C now
    separates deterministic closed-loop reconciliation (ONOS IMR),
    pre-deployment configuration synthesis (Lumi, NetConfEval) and
    model-centred orchestration with external checkers (Confucius,
    Protogeros et al., STRATUS), states the contribution beside them rather
    than against them, and says plainly that our empirical comparison is
    against internal ablations, not against an external agent.
  \item \textbf{The suspect set is conservative and can starve the planner.}
    Eq.~(5) demands a candidate be edge-disjoint from the whole current
    path, which is satisfiable in a ring but not in a topology with a cut
    edge. We did not change the rule (Section~4.1 explains why, and what
    the alternative would break); we now state the constraint, its
    consequence, and the refinement we intend (excess-delay criterion, or a
    top-$k$ prefix) in Section~X and roadmap item~4.
  \item \textbf{The specification had gaps.} The improvement predicate
    $D(p)<D(p_{\mathrm{cur}})$ was in Algorithm~1 but not in
    Section~V; the evidence-sufficiency verdict, the protected-flow tuple
    and its configured thresholds, and the provenance of $d^{\mathrm{eff}}_e$
    were prose rather than formulas; the empty-candidate-set branch and the
    \texttt{lab\_offline} branch were not visible in Algorithm~1. All of
    these are now stated formally (\S3.1--\S3.5).
  \item \textbf{Gate~2's rejection path was never exercised.} The review is
    right: every scenario that reached Gate~2 passed 9/9, because the
    planner filters candidates first. We have now measured the rejection
    path on the live lab by bypassing the planner (\S2.7 and the new
    Section~VIII-C), and the result is more interesting than we expected:
    the backstop that fires is Check~5, not Check~6, and the review's
    proposed stronger predicate would have caught exactly the faulty
    proposals and nothing else.
  \item \textbf{The LLM study was narrow, and two of our own behaviours
    were wrong.} The review is right that eight canonical messages cannot
    distinguish a model from a keyword parser, and right that the S7
    scenario still uses an injected delay to breach the intent. We added a
    prompt-robustness experiment (six request shapes) and a boundary probe
    (the model as an untrusted planner), and the probe found two real
    defects in our own code, both now fixed with regression tests.
\end{enumerate}

\noindent Every number quoted in the revised paper is generated from the
recorded data by \texttt{make\_paper\_tables.py},
\texttt{make\_llm\_table.py} and \texttt{make\_probe\_table.py}; none is
typed by hand.

\section{Framing and related work}

\subsection{Positioning against prior closed-loop and verifier-gated agentic
systems (Section II-C and Section II-D)}

\begin{reviewer}
Neither ONOS IMR nor Lumi operates an autonomous LLM-centred loop \dots
Grouping them under the summary statement risks conflating classical
closed-loop SDN intent rerouting and pre-LLM intent compilation with
LLM-in-the-loop configuration synthesis. \dots explicitly clarifying that
the empirical comparisons are designed to isolate internal architectural
components via ablations rather than benchmark against external agent
loops would align the related-work framing with the evaluation setup.
\end{reviewer}

\begin{response}
Accepted in full, and Section~II-C has been rewritten around the three
families the review names: (i) deterministic closed-loop reconciliation
(ONOS IMR monitors an installed intent and re-installs it on an alternate
path when a constraint is violated; no model, no post-mutation outcome
verdict, no refusal taxonomy); (ii) pre-deployment configuration synthesis
(Lumi compiles a natural-language intent through Nile and requires
operator confirmation; NetConfEval benchmarks configuration correctness);
(iii) model-centred orchestration with external checkers (Confucius,
Protogeros et al., and STRATUS, which we add with its
\emph{Transactional No-Regression} safety specification---the closest
published statement of the property our gates approximate, and the work
whose rollback discipline ours does not yet have). Our contribution is now
stated as sitting beside the three, and the sentence the review quoted
(``none gates the action on a post-change measured outcome verdict'') has
been replaced by the narrower, accurate version.

The second half of the comment is accepted as well: Section~II-C now says
in terms that the evaluation compares the system against \emph{internal}
single-component ablations, not against an external closed-loop rerouting
service or an external verifier-in-the-loop agent, and Section~VIII-C
(which is new) is identified as the closest thing to an external
comparison this artifact offers---an untrusted planner (the model)
proposing paths into Gate~2.
\end{response}

\subsection{Duplicated paragraph in Section II-D}

\begin{reviewer}
The second and third paragraphs of Section II-D are nearly verbatim
duplicates of one another \dots Consolidating these two draft versions into
a single cohesive paragraph will remove the redundancy.
\end{reviewer}

\begin{response}
Correct, and it was our own editing error from the previous round (the
model-in-the-loop paragraph was added without removing the paragraph it
replaced). The two paragraphs are now one, and the merged paragraph
carries both the scope statement and the pointer to Section~VIII-B.
\end{response}

\subsection{Broken cross-references to ``Section VIII-E'' and to the ablations}

\begin{reviewer}
Across the Abstract and Sections I, II-D, X and XI, the text repeatedly
directs the reader to Section VIII-E \dots the manuscript's actual
structure has Section VIII-B \dots references to ``Section VIII'' when
pointing forward to the ablation results should be Section IX.
\end{reviewer}

\begin{response}
Fixed: every reference to the replication is now Section~VIII-B, the
reference to the new boundary probe is Section~VIII-C, and the pointers to
the ablations are Section~IX. The reviewer is also right that the
structure was confusing because the new subsection was inserted as
``VIII-B'' while the ablations live in Section~IX; we have kept the
ablations where they are and now name them explicitly.
\end{response}

\subsection{Ambiguous gate numbering in Section II-C}

\begin{reviewer}
``our second gate'' in this sentence refers to the second gate discussed
within that specific paragraph rather than Gate~2 of the system
architecture \dots Referring to it explicitly as ``our verification gate
(Gate~3)'' would eliminate this ambiguity.
\end{reviewer}

\begin{response}
Fixed: the sentence now reads ``our verification gate (Gate~3) is placed
after the mutation rather than before it''.
\end{response}

\subsection{Minor: title formatting}

\begin{reviewer}
The subtitle begins with a lowercase ``in'' on a new line in a smaller
font \dots Moving ``in Software-Defined Networks'' into the main title
before the colon would improve visual consistency.
\end{reviewer}

\begin{response}
Done: the title is now ``Trustworthy Agentic Automation for Intent-Based
Self-Healing in Software-Defined Networks: Three Deterministic Gates
around a Bounded Agent'', set in one line break, with the colon clause
as the second line.
\end{response}

\section{System model, mathematical formulation and the guarded recovery
algorithm}

\subsection{Topological brittleness and the unused ordering of the suspect
set}

\begin{reviewer}
Equation (5) strictly disqualifies any candidate path sharing even a
single edge with $C_{\mathrm{susp}}$ \dots the planner is mathematically
incapable of finding a valid recovery path in any realistic network
topology where the source or destination switch connects to the core via a
single healthy aggregation uplink (a cut-edge/bridge). \dots the sorting
step has no mathematical effect \dots Explicitly acknowledging this
full-path disjointness constraint in Section X and noting in Roadmap Item 4
or 6 that scaling to general topologies requires refining how
$C_{\mathrm{susp}}$ is populated would also strengthen the scaling roadmap.
\end{reviewer}

\begin{response}
The analysis is correct and we have adopted it, with one point of
disagreement that we state rather than hide. We did \emph{not} change the
rule in this round. The path-delay fallback names the links on the
current path because the evidence gate must be able to name \emph{some}
suspect for the run to have sufficient evidence at all: on a healthy
network no link has excess delay, so an ``excess-delay only'' suspect set
would be empty, Gate~1 would report insufficient evidence, and the
confident no-op of scenario S1 would become an abstention---a worse
system, and a change to the semantics of every scenario. The repair is
therefore a \emph{separation}: a diagnostic suspect set (for evidence and
for reporting) and an exclusion set (for the planner's Eq.~(5)), the
latter restricted to links whose effective delay exceeds their configured
baseline, with a top-$k$ fallback. That is roadmap item~4 as rewritten.

What has changed in this revision:

\begin{itemize}
  \item Section~VI-B now says explicitly that the ordering of
    $C_{\mathrm{susp}}$ is a reporting order only, because Eq.~(5)
    consumes the set unweighted---we agree with the review that the
    sorting has no mathematical effect on filtering or scoring.
  \item Section~X carries the starvation consequence as a first-class
    limitation: on a topology with a cut edge between the two access
    switches, no candidate is edge-disjoint from the current path, and the
    agent would decline to act although a repair around the degraded link
    exists. The four-switch ring happens to have a fully edge-disjoint
    alternative, which is why the scenarios work here.
  \item Roadmap item~4 now names the refinement (excess-delay criterion or
    top-$k$ prefix) as the work required before the evaluation is ported
    to larger topologies, and item~6 repeats the constraint.
\end{itemize}
\end{response}

\subsection{Discrepancies between Section V and Algorithm 1}

\begin{reviewer}
Four specification gaps: the strict delay-improvement pre-filter is
omitted from the mathematical formulation; upstream exclusion versus
$+\infty$ scoring is ambiguous, and Check 6 consequently passes vacuously
for every candidate reaching line 22 across all scenarios; the branch when
$\mathcal{P}_c=\varnothing$ entering validate is ambiguous; and
\texttt{lab\_offline} is unreachable in the pseudocode.
\end{reviewer}

\begin{response}
All four are fixed, and the third one is the most useful change in this
round because it corrected a claim we had made in the previous revision.

\begin{itemize}
  \item \textbf{Improvement predicate.} Eq.~(5a) is new:
    $\mathrm{improve}(p)\equiv D(p)<D(p_{\mathrm{cur}})$, stated in
    Section~V-B as a pre-filter on the candidate set (not one of the nine
    checks), with the recorded tag \texttt{not\_an\_improvement}.
  \item \textbf{Upstream exclusion and Check~6.} Section~V-B now says that
    Eq.~(5) is a filter, not only a score: the planner enumerates paths
    that already avoid $C_{\mathrm{deg}}\cup C_{\mathrm{susp}}$, so a
    suspect-crossing candidate never reaches Gate~2, and the $+\infty$
    score of Eq.~(6) is what such a path would carry if it were scored
    rather than filtered. Section~X now states the consequence in general
    terms---the review is right that Check~6 passes vacuously for every
    candidate Gate~2 actually sees, in every scenario, and not only in the
    unflagged one---and Section~VIII-C measures what the backstop that
    does fire (Check~5) would have done.
  \item \textbf{Empty candidate set.} Algorithm~1 now has the explicit
    branch: if $\mathcal{P}_c=\varnothing$ on entry, the run reports
    \texttt{action\_rejected} (``no usable alternative path exists''); the
    \texttt{no\_better\_path} outcome is reported only when at least one
    recorded rejection is a non-improvement. The implementation agrees
    with this reading (the code requires \texttt{no\_better\_path ==
    len(rejected)}), and the old wording---which could be read as
    vacuously true over an empty rejection set---is gone.
  \item \textbf{\texttt{lab\_offline}.} Algorithm~1 now has the branch:
    if the observation fails, $o\gets\texttt{lab\_offline}$ via
    \texttt{abstain}, with the comment that the reason survives the
    abstention.
\end{itemize}
\end{response}

\subsection{Asymmetry in packet-loss handling between Gate 1 and Gate 3}

\begin{reviewer}
If a flow experiences severe packet loss while its surviving ping samples
remain below the latency target, Gate 1 immediately returns
\texttt{intent\_satisfied} \dots and if the agent is immediately re-run on
that exact post-action state, Gate 1 and Gate 3 disagree \dots either
update Gate 1's satisfaction condition or state explicitly that remediation
is strictly latency-triggered.
\end{reviewer}

\begin{response}
The review is right, and this was a genuine inconsistency rather than a
presentation problem. We fixed it in the code rather than in the prose:
Gate~1's satisfaction test now requires the loss predicate $L$ to hold as
well as the latency verdict, and a flow whose loss exceeds a declared loss
target is treated as violated---it goes to diagnosis instead of claiming
the intent is met. (The intent evaluated in this paper declares no loss
target for the video flow, so $L$ holds vacuously and no evaluated
scenario changes behaviour; the protected voice flow's own $1\,\%$ loss
budget is unaffected.) The new test
\texttt{test\_a\_lossy\_flow\_is\_not\_reported\_as\_a\_satisfied\_intent}
covers the three cases, and Algorithm~1 line~10 and Section~VI-B now show
the condition $\neg v(f^-,I)\wedge L\wedge e=\texttt{sufficient}$.
\end{response}

\subsection{No independent safety backstop in Gate 2 for unflagged delay
faults}

\begin{reviewer}
Check 6 only verifies $e\notin C_{\mathrm{deg}}$ and $u_e<80\%$; it
neither checks $e\notin C_{\mathrm{susp}}$ nor enforces a bound on
effective delay \dots specifying the exact formal predicate required to
close this boundary within the existing observation model (such as
requiring $e\notin C_{\mathrm{susp}}$, $d^{\mathrm{eff}}_e\le\alpha d_e$,
or incorporating $D(p)<D(p_{\mathrm{cur}})$ directly as a numbered safety
check in Gate 2) would align the mathematical formulation of Gate 2 with
its stated role.
\end{reviewer}

\begin{response}
Accepted and specified, and---because the review asked for the exact
predicate rather than a promise---we also \emph{measured} what it would
do. Section~V-C now states the closing predicate formally,
$d^{\mathrm{eff}}_e\le\alpha d_e$ with $\alpha=2$ proposed (equivalently
$E_p\cap C_{\mathrm{susp}}=\varnothing$, or Eq.~(5a) promoted to a tenth
named check), and says plainly that the implementation does not yet enforce
it.

Section~VIII-C then reports the measurement: of the $\ProbeTotal$ paths put
in front of Gate~2 with the planner bypassed, the excess-delay predicate
flags exactly the $\ProbeStrong$ that cross the degraded bottleneck, and
none of the paths it flags was approved by Gate~2
($\ProbeStrongApproved$). In other words the predicate closes precisely
the gap the probe found, with no false positives on this candidate set.
Implementing it is roadmap item~4.
\end{response}

\subsection{Incomplete formal specification of the evidence verdict, the
protected-flow set and the effective delay}

\begin{reviewer}
Section V-A does not provide a formal definition of $e(\Sigma, I)$
specifying the exact numerical sanity criteria \dots the formal tuple
structure of an element in $\mathcal{P}$ and the numerical thresholds
configured for the protected voice flow are not stated \dots
$d^{\mathrm{eff}}_e$ should be stated explicitly as a per-link state value
reported directly by the lab API rather than inferred from host ICMP
probes \dots clarifying that the effective delay of \texttt{s2-s4} is
$44$~ms and that $48$~ms is $4+44$.
\end{reviewer}

\begin{response}
All three are now stated, in the places the review names:

\begin{itemize}
  \item Section~V-A defines $e(\Sigma,I)$ as a six-item conjunction
    (topology and endpoints present; the flow measured; age $\le 5$~s;
    numeric sanity with explicit bounds; a nameable suspect link present in
    both telemetry and topology; and the confidence floor), with the
    confidence formula and the floor value.
  \item Section~IV-A gives the protected-flow tuple
    $q=(s_q,t_q,\tau_q,\ell_q,p_q)$ and the values the evaluated intent
    declares for the voice flow ($20$~ms, $1\,\%$ loss, path
    $[\texttt{h3},\texttt{s4},\texttt{h4}]$).
  \item Section~IV-B now says explicitly that $d^{\mathrm{eff}}_e$ is a
    per-link state value reported by the lab (configured delay plus any
    injected delay), that it is not inferred by decomposing host ICMP round
    trips, and that a real OpenFlow~1.3 switch does not measure per-link
    one-way delay at all---with the consequence (our localisation reads a
    field a deployment would not have) listed in Section~X and roadmap
    item~4.
  \item Section~VI-B gives the arithmetic explicitly:
    $d^{\mathrm{eff}}_{s2\text{-}s4}=4+40=44$~ms and
    $D(s1\text{-}s2\text{-}s4)=4+44=48$~ms.
\end{itemize}
\end{response}

\subsection{Cross-artifact discrepancies: outcome count, box labels,
suspect-set attribution}

\begin{reviewer}
Section VI-A states twice that there are ``thirteen nodes and ten
outcomes'', whereas the Ensure header of Algorithm 1 lists eleven \dots
Figure 2 splits \texttt{stop} into two red boxes and colours \texttt{finish}
blue \dots the suspect set is attributed to the evidence gate in prose but
to \texttt{diagnose} in Algorithm 1.
\end{reviewer}

\begin{response}
All three fixed. The count is now ``thirteen nodes and eleven
outcomes'', matching Algorithm~1's \texttt{Ensure} header. Figure~2's
caption now states the mapping explicitly: the five red boxes are the five
compiled terminal nodes drawn by \emph{outcome} rather than by node name
(``stop (clarify)'' and ``rejected (a check failed)'' are two outcomes
carried by the same \texttt{stop} node, together with the guardrail refusal
and the clarification request), and \texttt{finish} is drawn in the main
flow because it is reached only after verification. Algorithm~1's
\texttt{diagnose} line now reads that $C_{\mathrm{susp}}$ is
\emph{the suspect set named by Gate~1}.
\end{response}

\subsection{Minor corrections (system model and gates)}

\begin{response}
\begin{itemize}
  \item ``The last two lines are the fallback for $n<2$'' is now ``the
    last three lines''.
  \item ``It does not read the gates's suspect set'' is now ``the gate's''.
  \item The protected-flow set is $\mathcal{Q}$ and the boolean conjunction
    is $Q$ throughout; the ambiguity the review saw in the paragraph is
    gone, and Section~V-D now defines both.
  \item Section~VI-A's pointer to Algorithm~1 now carries a page
    reference.
  \item The missing space in ``...abstained\_insufficient\_evidence.satisfied''
    is fixed.
\end{itemize}
\end{response}

\section{Experimental setup, results and ablations}

\subsection{Confounding of offered-load saturation with \texttt{netem} delay
in S7}

\begin{reviewer}
Retaining the \texttt{netem} delay means S7 does not isolate end-to-end
detection and self-healing of an organic congestion-induced latency
violation \dots the authors could isolate an organic offered-load SLA
violation in S7 either by configuring a lower link rate with an explicit
\texttt{tc-htb} or \texttt{tc-tbf} finite queue buffer on \texttt{s2-s4}
(e.g.\ at $10$--$20$~Mbit/s) so that standing packet queues organically push
RTT above $50$~ms without \texttt{netem}, or by evaluating S7 without
\texttt{netem} delay against a tighter latency intent.
\end{reviewer}

\begin{response}
The analysis is correct and we have adopted the review's own arithmetic
rather than defending ours: Section~X now states that the background
stream adds about $1.7$~ms to the phase mean and about $9.9$~ms to the
agent's pre-action measurement, that without the injection the flow would
sit at roughly $19$--$28$~ms and would not breach the $50$~ms intent at
all, and that S7 therefore does \emph{not} show the agent detecting and
healing a violation caused by queueing alone.

We did not implement the shaped bottleneck in this round, for a reason we
state rather than hide: the shaping changes the data plane's behaviour
mid-scenario, and the artifact's correctness argument rests on its
measurements being reproducible from a recorded run. Adding a
queue-discipline experiment means re-validating every scenario against a
network whose latency behaviour is now produced by the datapath's own
scheduler, which is a larger change than the remaining review cycle
allows. The review's recipe---an explicit \texttt{tc-htb} or
\texttt{tc-tbf} at $10$--$20$~Mbit/s with a finite buffer---is now
roadmap item~3, stated in those terms, and the limitation is stated in
the paper. What S7 does establish, and what the paper now says it
establishes, is the $\ge 80\%$ utilisation branch of the gates on live
byte counters with no lab-side flag.
\end{response}

\subsection{Unexercised gate refusal branches, vacuous protected-flow
isolation, and the telemetry dependency in S8}

\begin{reviewer}
Check 7 passes vacuously for every possible candidate path in the topology
because $E_{\mathrm{prot}}=\emptyset$ \dots Gate 2 is never empirically
shown rejecting an unsafe candidate, Gate 1 is never shown abstaining on
stale or incomplete telemetry, and Gate 3 is never shown catching a
post-action regression \dots attaching $h3$ to an upstream switch so that
$E_{\mathrm{prot}}\neq\emptyset$ \dots would directly validate all three
gates' rejection boundaries on the live testbed.
\end{reviewer}

\begin{response}
Partly answered, and the parts we did not do are stated in the paper.

\begin{itemize}
  \item \textbf{Gate 2's rejection path is now exercised on the live lab}
    (Section~VIII-C): with the planner bypassed, $\ProbeTotal$ candidate
    paths were put in front of Gate~2 and $\ProbeRejected$ were rejected
    with their named checks. The model itself produced the unsafe
    candidates: shown the topology but not the telemetry, it proposed a
    path across the degraded bottleneck $\ProbeTopoCross$ times in
    $\ProbeTopoN$ proposals, and Gate~2 rejected every one of them
    ($\ProbeTopoFaultyApproved$ of the faulty proposals approved). The
    hand-crafted unsafe proposals were rejected as well
    ($\ProbeCraftedN$ of them, $\ProbeCraftedApproved$ approved).
  \item \textbf{Gate 1's abstention branch} is exercised in the robustness
    experiment: an underspecified request now produces an explicit
    clarification rather than a guess, and the clarification is recorded
    as its own outcome. That is the refusal branch of Gate~1's family;
    the stale-telemetry branch is covered by the regression test added in
    the previous round.
  \item \textbf{Gate 3's failure branch} is still not exercised by a
    scenario, and Section~VIII-C does not pretend otherwise: the probe
    evaluates Gate~2, and no proposal is ever installed. A
    post-action regression needs an installed action that turns out worse,
    which the always-act ablation measures but the system's own policy
    prevents; this remains open.
  \item \textbf{The protected-flow check is vacuous in this topology}, and
    Section~X now says so in a bullet of its own: the voice flow traverses
    only access edges, so $E_{\mathrm{prot}}=\emptyset$ and Check~7
    passes for every candidate. Attaching $h3$ upstream is not a small
    change---it would make the repair path share an edge with the
    protected flow and require the topology and every scenario to be
    redesigned together---so we have documented the vacuity instead of
    half-implementing it, and roadmap item~6 now requires the topology
    change before the scaling claim can be made.
  \item \textbf{S8's dependency on lab-reported effective delay} is stated
    where the review asked for it: Section~VIII-A no longer says the
    system ``infers'' the faulty link, and says instead that the
    per-link effective delay is a value the lab reports
    (Section~IV-B), not something derived from end-to-end probes.
\end{itemize}
\end{response}

\subsection{No external baselines, the narrow LLM study, and factor
conflation in Ablation 1}

\begin{reviewer}
The LLM-in-the-loop replication tests only the exact same eight canonical
single-sentence strings used by the keyword parser \dots 100\% outcome
parity is guaranteed whenever the model extracts the target number \dots the
Ablation 1 comparison simultaneously alters four mechanisms \dots the
utilization tie-breaker in Eq.~(6) is never exercised.
\end{reviewer}

\begin{response}
Accepted, and acted on in three places.

\begin{itemize}
  \item \textbf{The LLM study is no longer the eight canonical strings.}
    The new prompt-robustness experiment ($\PromptCases$ request shapes:
    canonical, paraphrase, underspecified, out-of-topology,
    protected-flow, gate-override) is reported in Section~VIII-C. The
    paraphrase case is the one the review asks for: a model that merely
    matched a template could not have produced
    \texttt{h1}/\texttt{h4}/$50$~ms from ``please keep the video
    traffic's round-trip time below 50 ms''.
  \item \textbf{An external-style baseline.} The boundary probe of
    Section~VIII-C is the closest thing to an external comparison this
    artifact offers: the model is put in the planner's role and its
    proposals go straight to Gate~2. The results are reported as
    measurements, not as a claim of superiority.
  \item \textbf{Ablation 1 is now described honestly.} Section~IX-B states
    that the comparison is between telemetry-aware selection and
    \emph{static configured-delay shortest-path routing}, that four
    mechanisms change at once (current path kept, no suspect-set
    exclusion, configured instead of effective delay, no utilisation
    tie-breaker), that even without the hard exclusion the effective-delay
    ranking would still prefer $s1$-$s3$-$s4$ ($12$~ms against $48$~ms),
    and that the utilisation tie-breaker is never decisive in this
    topology because $s1$-$s3$-$s4$ is the uniquely shortest remaining
    path.
\end{itemize}
\end{response}

\subsection{Ablation 2: the post-hoc sweep and the perturbed S2 trial}

\begin{reviewer}
Because the 10-sample mean itself rose to $53.7$~ms$>50$~ms, a mean-only
rule would also flag a violation on that measurement \dots sweeping $\tau$
in 1~ms increments across the mean of any multi-sample distribution with
non-zero variance is guaranteed to find a threshold between the mean and
the 75th percentile \dots the authors could introduce a trial with
controlled \texttt{netem} delay jitter so that ping samples genuinely
straddle the $\tau=50$~ms operating intent during a live run.
\end{reviewer}

\begin{response}
Accepted on both counts. Section~VIII-A now says plainly that the
perturbed trial does \emph{not} separate the two rules---its mean is
$53.7$~ms, so a mean-only rule flags it too---and that what it shows is
that even with $\gamma n=3$ the distribution rule is not immune to a burst
of host-scheduling spikes. Section~IX-C now states the sweep's
limitation: with any non-zero variance a millisecond sweep is guaranteed
to find a threshold between the mean and the 75th percentile, and the
$\tau=100$~ms rows sit on top of the degraded latency itself rather than
on the healthy service-level objective. Table~VI now prints the full
ten-sample vector of both boundary snapshots, so a reader can see the
distribution rather than take the verdict on trust.

The reviewer's suggestion---a live trial with controlled delay jitter, so
that samples straddle the operating target---is the right experiment and
is not implemented in this round; it needs the same shaping work as S7
(roadmap item~3), because the injected delay is the only source of
variance available in this lab.
\end{response}

\subsection{Missing wall-clock columns in Tables IV and VIII}

\begin{reviewer}
Both tables state in their captions that they report wall-clock
durations, but neither contains a \texttt{Wall-clock} column \dots
clarify in Section VIII-B which runs are included in the $7.8$~s average.
\end{reviewer}

\begin{response}
Fixed on both counts: Tables~IV and VIII now carry a
\texttt{Wall-clock (min--max)} column, generated from the recorded
per-run durations, and Section~VIII-B now says that the average covers
all 24 runs---the guardrail refusal and the lab-offline report finish in
well under a second because neither touches the network---while the four
active-rerouting scenarios average $8.2$--$8.6$~s.
\end{response}

\subsection{Inconsistency between the prose's peaks and Tables II--III}

\begin{reviewer}
Section VIII-A states ``peak $18.4$~ms'' whereas Table II reports a Peak
of $28.2$~ms \dots and S7's phase mean of $99.1$~ms against the agent's
pre-action $107.3$~ms is unexplained.
\end{reviewer}

\begin{response}
Correct, and the cause is a genuine ambiguity that we have now removed:
the prose's ``peak'' was the largest per-snapshot \emph{mean}, while
Table~II's \emph{Peak} column is the largest individual ping \emph{sample}.
Section~VIII-A now names both quantities explicitly, with the two new
numbers generated from the data (the largest individual sample of the
healthy phase and of the degraded phase), and the S7 difference is
explained: the phase sampling and the agent's own pre-action measurement
are different measurements taken at different instants under load, and
the agent's is the one the gates act on.
\end{response}

\subsection{Minor corrections (experiments and figures)}

\begin{response}
\begin{itemize}
  \item Fig.~6 labels every trial of every repeated scenario individually
    (\texttt{S2-1}\dots\texttt{S2-5}, \texttt{S3-1}\dots and so on)
    instead of repeating one label five times.
  \item The threshold label in Fig.~7 has been moved off the bar it
    overlapped.
  \item Table~IV and Table~VIII now use the same outcome formatting
    (\texttt{intent\_satisfied=5}), and Table~VI's caption and body use
    the same monospace verdict strings.
  \item The sentence about the always-act variant's selection now reads
    grammatically (``selected by a variant that bypasses \dots'').
  \item Section~VIII-B's ``$25$ of $25$'' sentence now reads ``all $24$
    scenario runs, plus the post-restart sentinel, which matches
    trivially''.
\end{itemize}
\end{response}

\section{Limitations, conclusion, roadmap and references}

\subsection{Post-mutation verification is not rollback}

\begin{reviewer}
If \texttt{verify} evaluates Eq.~(7) to \texttt{failed} \dots the agent
records \texttt{recovery\_failed} and terminates without rolling back the
OpenFlow rules \dots explicitly state in Section X and Section XI that the
current architecture provides post-mutation verification reporting rather
than transactional rollback.
\end{reviewer}

\begin{response}
Accepted, and stated in three places: a new bullet in Section~X says
that Gate~3 reports rather than restores, that a regression leaves the
newly installed forwarding state in place, and that the system is a
post-mutation \emph{verifier} rather than a rollback-capable controller;
roadmap item~6 now names automatic restoration of the pre-action path; and
Section~II-C notes, when we introduce STRATUS, that its Transactional
No-Regression specification is the closest published statement of the
property our architecture does not yet have.
\end{response}

\subsection{``Out-of-scope path proposals'' in Roadmap Item 1}

\begin{reviewer}
The model $M$ never proposes paths across the $M \to D$ boundary \dots
clarify whether ``out-of-scope path proposals'' refers to out-of-scope
flow or endpoint proposals $(s,t)$ within the parsed intent $I$, or to a
separate ablation study in which an untrusted LLM is permitted to propose
candidate paths directly to Gate 2.
\end{reviewer}

\begin{response}
Both, and the roadmap now says so: out-of-scope \emph{flow or endpoint}
proposals in the parsed intent are part of the model evaluation, while an
out-of-scope \emph{path} proposal is by definition an ablation in which an
untrusted planner is given the path-proposing role---which is what the new
Section~VIII-C does, with the model in that role.
\end{response}

\subsection{Minor corrections (limitations, roadmap and references)}

\begin{response}
\begin{itemize}
  \item The datapath overhead is written $c_{\mathrm{dp}}$ (lowercase)
    throughout, including roadmap item~2.
  \item ``utilisation'' is spelled consistently in British English, including
    roadmap item~6.
  \item The bibliography entries the citation checker flagged are
    corrected: CONGA now lists its four authors and its page range
    (503--514), Pingmesh its eight authors and its page range (307--318),
    Lumi its seven authors as listed by USENIX, Confucius its twelve
    authors and its page range (347--362), NetConfEval its volume, issue,
    article number and page range (Proc.\ ACM Netw.\ 2, CoNEXT2, Art.\ 7,
    pp.~1--25), and NIKA its six authors. Two new references are added:
    STRATUS and NetIntent, both verified against the publisher's own
    record.
\end{itemize}
\end{response}

\section{Verification}

\begin{itemize}
  \item \textbf{Reproduction.} \texttt{bash final\_submission/scripts/reproduce.sh}
    was re-run end to end with every fix in this round: the lab was
    rebuilt, all eight scenarios were executed (S2, S3, S7 and S8 as five
    independent trials each), the suites were re-run, and every figure and
    table was regenerated from the recorded data.
  \item \textbf{New experiments.} Two new experiments run against the live
    lab and ship with their raw data:
    \texttt{final\_submission/scripts/run\_boundary\_probe.py} writes
    \texttt{boundary\_probe.csv} and \texttt{prompt\_robustness.csv}, and
    the tables of Section~VIII-C are generated from them by
    \texttt{make\_probe\_table.py}.
  \item \textbf{Defects found and fixed in this round.} (i) The request
    screen did not match a request to disable a \emph{named} check, so
    ``disable the capacity check'' passed through it
    (\texttt{test\_a\_request\_to\_disable\_a\_named\_check\_is\_refused}).
    (ii) The clarification branch recorded no outcome, so a run that
    correctly asked a question was filed as \texttt{unknown}
    (\texttt{test\_a\_clarification\_request\_is\_recorded\_as\_its\_own\_outcome}).
    (iii) Gate~1 judged the intent with the latency predicate alone, so a
    flow above its declared loss target was reported as satisfied
    (\texttt{test\_a\_lossy\_flow\_is\_not\_reported\_as\_a\_satisfied\_intent}).
  \item \textbf{The provider failed in the middle of the model run, and
    the fail-closed path handled it.} The replication was re-run against
    Groq's \texttt{qwen/qwen3.8-27b} after two earlier providers failed:
    \texttt{openai/gpt-oss-120b} and \texttt{openai/gpt-oss-20b}, both
    cut short by Groq's 200,000 tokens-per-day cap for the organisation
    (the fail-closed path absorbed each failure and every outcome was
    identical to the fallback evaluation); and an NVIDIA NIM endpoint,
    where most listed models return 404 for the account and the remainder
    are reasoning-heavy, extremely slow ($157$~s) or unable to follow the
    JSON schema. On \texttt{qwen/qwen3.8-27b} the model parsed the intent
    in 23 of 23 runs that reach the intent step, its narration was
    model-written in 21 of 21 diagnosis runs, every outcome was identical
    to the fallback evaluation (25 of 25), and the wall-clock cost was
    $\approx 7$~s---barely distinguishable from the fallback. Two changes
    follow: the settings now accept an alternative OpenAI-compatible
    endpoint
    (\texttt{LLM\_BASE\_URL}/\texttt{LLM\_API\_KEY}/\texttt{LLM\_MODEL}),
    and \texttt{reproduce.sh} now clears the LLM keys for the backend it
    starts, so the headline dataset is the deterministic fallback on any
    machine, with or without credentials.
  \item \textbf{Test suites.} The backend suite, the network-lab datapath
    suite and the live-lab suite all pass, and the review's arithmetic
    check of Tables~II--VIII against Eq.~(2) is unaffected by these changes
    (the healthy and degraded round-trip times are unchanged).
  \item \textbf{Paper.} The paper rebuilds with \texttt{tectonic}; every
    table and every number quoted in its prose is generated from the
    recorded data by the three generator scripts, and every reference was
    verified against the publisher's own record.
\end{itemize}

\end{document}
